\documentclass[11pt]{article}
\usepackage{amsmath,amssymb,amsfonts}
\usepackage{geometry}
\usepackage{algorithm}
\usepackage{algpseudocode}
\usepackage{hyperref}
\geometry{margin=1in}

\title{Pixel- and Image-Based Pathfinding Algorithms:\\
Design, Pseudocode, and Incremental Enhancements}
\author{Shaimaa Said Soltan}
\date{June 1, 2026}

\begin{document}
\maketitle

\tableofcontents

\section{Introduction}

This document describes a family of pathfinding and traversal algorithms
designed for road networks and maze-like images, all implemented in a
pixel- or grid-based setting. Each section presents:
\begin{itemize}
  \item The core idea of the algorithm.
  \item Detailed pseudocode.
  \item Enhancements and how it advances beyond previous variants.
\end{itemize}

We assume a 2D discrete domain (image or rasterized road network) with
width $W$ and height $H$, indexed by integer coordinates
$(x,y)$, $0 \le x < W$, $0 \le y < H$.

%%%%%%%%%%%%%%%%%%%%%%%%%%%%%%%%%%%%%%%%%%%%%%%%%%%%%%%%%%%%
\section{Pixel-Based Maze Solver}
\label{sec:pixel_maze}

\subsection{Concept}

The pixel-based maze solver treats the rasterized road network (or maze)
as a binary image: \emph{road} pixels (white) and \emph{non-road} pixels
(black). A standard breadth-first search (BFS) is performed on the
4-neighborhood (up, down, left, right) of each pixel to find a shortest
path between a start pixel $s$ and a target pixel $t$.

\subsection{Pseudocode}

We assume a function \texttt{isRoad(x,y)} that returns \texttt{true} if
pixel $(x,y)$ is a valid road pixel.

\begin{algorithm}[h!]
\caption{Pixel-Based Maze Solver (4-Connected BFS)}
\begin{algorithmic}[1]
\Require Start pixel $s = (s_x,s_y)$, target pixel $t = (t_x,t_y)$
\Require Image dimensions $W,H$
\Ensure Path $P$ as a sequence of pixels from $s$ to $t$, or empty if none

\State Initialize \texttt{visited[0..W*H-1]} to \texttt{false}
\State Initialize \texttt{prev[0..W*H-1]} to \texttt{-1}
\State Initialize queue $Q$ as empty

\State $idx_s \gets s_y \cdot W + s_x$
\State \texttt{visited[$idx_s$]} $\gets$ \texttt{true}
\State Enqueue $s$ into $Q$

\State Define neighbor offsets:
\State \hspace{1em} $\Delta = \{(1,0),(-1,0),(0,1),(0,-1)\}$

\State $found \gets false$

\While{$Q$ is not empty}
  \State $(x,y) \gets$ Dequeue from $Q$
  \If{$(x,y) = (t_x,t_y)$}
    \State $found \gets true$
    \State \textbf{break}
  \EndIf
  \ForAll{$(dx,dy) \in \Delta$}
    \State $nx \gets x + dx$, $ny \gets y + dy$
    \If{$0 \le nx < W$ and $0 \le ny < H$}
      \State $idx \gets ny \cdot W + nx$
      \If{\texttt{visited[$idx$]} = \texttt{false} and \texttt{isRoad(nx,ny)}}
        \State \texttt{visited[$idx$]} $\gets$ \texttt{true}
        \State \texttt{prev[$idx$]} $\gets y \cdot W + x$
        \State Enqueue $(nx,ny)$ into $Q$
      \EndIf
    \EndIf
  \EndFor
\EndWhile

\If{\textbf{not} $found$}
  \State \textbf{return} empty path
\EndIf

\State // Reconstruct path
\State $P \gets$ empty list
\State $idx \gets t_y \cdot W + t_x$
\While{$idx \neq -1$}
  \State $x \gets idx \bmod W$
  \State $y \gets \lfloor idx / W \rfloor$
  \State Prepend $(x,y)$ to $P$
  \State $idx \gets \texttt{prev}[idx]$
\EndWhile

\State \textbf{return} $P$
\end{algorithmic}
\end{algorithm}

\subsection{Enhancements and Advancement}

\begin{itemize}
  \item \textbf{Enhancement:} Using a rasterized representation of roads
  allows direct BFS on pixels, avoiding explicit graph construction.
  \item \textbf{Advancement:} This algorithm is the baseline for all
  subsequent methods. Later algorithms extend it with more directions,
  different neighbor models, or additional constraints (e.g., longest
  path, directional bias).
\end{itemize}
\clearpage
%%%%%%%%%%%%%%%%%%%%%%%%%%%%%%%%%%%%%%%%%%%%%%%%%%%%%%%%%%%%
\section{Longest-White-Path Solver}
\label{sec:longest_white}

\subsection{Concept}

Instead of finding the \emph{shortest} path between two points, the
longest-white-path solver aims to find a path that maximizes the number
of white (road) pixels visited, subject to connectivity constraints.
This is generally a hard problem (longest path in a graph is NP-hard),
so we use a heuristic or constrained variant.

A simple variant: perform a BFS/DFS from a start pixel and track the
longest simple path discovered under a depth or time limit.

\subsection{Pseudocode (Heuristic DFS-Based Longest Path)}

\begin{algorithm}[h!]
\caption{Heuristic Longest-White-Path Solver}
\begin{algorithmic}[1]
\Require Start pixel $s = (s_x,s_y)$
\Require Image dimensions $W,H$
\Ensure A long (not necessarily optimal) path $P$

\State $bestPath \gets$ empty list

\State Define neighbor offsets:
\State \hspace{1em} $\Delta = \{(1,0),(-1,0),(0,1),(0,-1)\}$

\State Initialize \texttt{visited[0..W*H-1]} to \texttt{false}

\Function{DFS}{$x,y,currentPath$}
    \State Append (x,y) to currentPath
    \State $idx \gets y \cdot W + x$
    \State visited[idx] $\gets$ true

  \If{$|currentPath| > |bestPath|$}
    \State $bestPath \gets currentPath$
  \EndIf

  \ForAll{$(dx,dy) \in \Delta$}
    \State $nx \gets x + dx$, $ny \gets y + dy$
    \If{$0 \le nx < W$ and $0 \le ny < H$}
      \State $nidx \gets ny \cdot W + nx$
      \If{\texttt{visited[$nidx$]} = \texttt{false} and \texttt{isRoad(nx,ny)}}
        \State DFS($nx,ny$, $currentPath$)
      \EndIf
    \EndIf
  \EndFor

  \State \texttt{visited[$idx$]} $\gets$ \texttt{false}
  \State Remove last element from $currentPath$
\EndFunction

\State DFS($s_x,s_y$, empty list)
\State \textbf{return} $bestPath$
\end{algorithmic}
\end{algorithm}

\subsection{Enhancements and Advancement}

\begin{itemize}
  \item \textbf{Enhancement:} Instead of minimizing path length, this
  solver explores the space of simple paths to maximize coverage of
  white pixels.
  \item \textbf{Advancement:} This algorithm extends the pixel-based
  maze solver by changing the objective from shortest path to longest
  path, illustrating how the same underlying grid representation can
  support different optimization goals.
\end{itemize}
\clearpage
%%%%%%%%%%%%%%%%%%%%%%%%%%%%%%%%%%%%%%%%%%%%%%%%%%%%%%%%%%%%
\section{16-Direction BFS Solver}
\label{sec:16_dir_bfs}

\subsection{Concept}

The 4-connected BFS only allows moves in cardinal directions. The
16-direction BFS solver extends the neighbor set to include diagonals
and additional intermediate directions, approximating a more continuous
movement model and reducing path jaggedness.

We define a set of 16 direction vectors $(dx_i,dy_i)$, each representing
a possible step.

\subsection{Pseudocode}

\begin{algorithm}[h!]
\caption{16-Direction BFS Solver}
\begin{algorithmic}[1]
\Require Start pixel $s$, target pixel $t$, image dimensions $W,H$
\Ensure Shortest path in terms of number of steps (in 16-connected grid)

\State Initialize \texttt{visited}, \texttt{prev}, and queue $Q$ as in
Section~\ref{sec:pixel_maze}.

\State Define 16-direction neighbor offsets:
\State \hspace{1em} $\Delta \gets \{(1,0),(-1,0),(0,1),(0,-1),$
\State \hspace{4em} $(1,1),(1,-1),(-1,1),(-1,-1),$
\State \hspace{4em} $(2,1),(2,-1),(-2,1),(-2,-1),$
\State \hspace{4em} $(1,2),(1,-2),(-1,2),(-1,-2)\}$

\State Run BFS exactly as in Algorithm~1, but using $\Delta$ above.

\State Reconstruct path using \texttt{prev} as before.
\end{algorithmic}
\end{algorithm}

\subsection{Enhancements and Advancement}

\begin{itemize}
  \item \textbf{Enhancement:} The neighbor set is expanded from 4 to 16
  directions, allowing more natural diagonal and shallow-angle moves.
  \item \textbf{Advancement:} Compared to the pixel-based maze solver,
  this algorithm produces smoother, less Manhattan-like paths and better
  approximates Euclidean shortest paths on the raster.
\end{itemize}
\clearpage
%%%%%%%%%%%%%%%%%%%%%%%%%%%%%%%%%%%%%%%%%%%%%%%%%%%%%%%%%%%%
\section{Grid-Based Neighbor Solver}
\label{sec:grid_neighbor}

\subsection{Concept}

The grid-based neighbor solver abstracts the image into a coarser grid
(e.g., superpixels or tiles). Each cell in the coarse grid represents a
block of pixels, and neighbors are defined at the cell level. This
reduces the search space and can accelerate pathfinding on large images.

\subsection{Pseudocode}

Assume we partition the image into cells of size $k \times k$ pixels.
Let $(i,j)$ index cells, with $0 \le i < W_c$, $0 \le j < H_c$.

\begin{algorithm}[h!]
\caption{Grid-Based Neighbor Solver}
\begin{algorithmic}[1]
\Require Coarse grid dimensions $W_c,H_c$, cell size $k$
\Require Start cell $S_c$, target cell $T_c$
\Ensure Path as a sequence of cells (optionally refined to pixels)

\State Define \texttt{cellIsRoad(i,j)} = \texttt{true} if the cell
contains at least one road pixel.

\State Initialize \texttt{visited[0..W$_c$H$_c$-1]} and
\texttt{prev[0..W$_c$H$_c$-1]} as before.

\State Define neighbor offsets on the coarse grid:
\State \hspace{1em} $\Delta_c = \{(1,0),(-1,0),(0,1),(0,-1)\}$

\State Run BFS on cells $(i,j)$ using \texttt{cellIsRoad} and $\Delta_c$.

\State Reconstruct cell-level path $P_c$ from $S_c$ to $T_c$.

\State Optionally, for each consecutive pair of cells in $P_c$, refine
to a pixel-level path inside the union of the two cells using a local
pixel-based BFS.

\State \textbf{return} refined path.
\end{algorithmic}
\end{algorithm}

\subsection{Enhancements and Advancement}

\begin{itemize}
  \item \textbf{Enhancement:} Coarse grid representation reduces the
  state space, improving performance on large maps.
  \item \textbf{Advancement:} This algorithm generalizes the pixel-based
  solver by introducing a multi-resolution view: global routing on a
  coarse grid, local refinement on pixels.
\end{itemize}
\clearpage
%%%%%%%%%%%%%%%%%%%%%%%%%%%%%%%%%%%%%%%%%%%%%%%%%%%%%%%%%%%%
\section{Image-Based Traversal Solver}
\label{sec:image_traversal}

\subsection{Concept}

The image-based traversal solver treats the entire image as a weighted
graph, where each pixel has a traversal cost derived from image
properties (e.g., brightness, color, road thickness). Instead of a
binary \texttt{isRoad}, we use a cost function $c(x,y)$ and run a
Dijkstra-like algorithm or weighted BFS.

\subsection{Pseudocode (Dijkstra on Image)}

\begin{algorithm}[h!]
\caption{Image-Based Traversal Solver (Dijkstra)}
\begin{algorithmic}[1]
\Require Start pixel $s$, target pixel $t$, cost function $c(x,y)$
\Ensure Minimum-cost path from $s$ to $t$

\State Initialize distance array \texttt{dist[0..W*H-1]} to $+\infty$
\State Initialize \texttt{prev[0..W*H-1]} to \texttt{-1}
\State Use a priority queue $PQ$ keyed by distance

\State $idx_s \gets s_y \cdot W + s_x$
\State \texttt{dist[$idx_s$]} $\gets 0$
\State Insert $(0, s_x, s_y)$ into $PQ$

\State Define neighbor offsets (4 or 8 directions):
\State \hspace{1em} $\Delta = \{(1,0),(-1,0),(0,1),(0,-1)\}$

\While{$PQ$ not empty}
  \State Extract $(d,x,y)$ with minimal $d$ from $PQ$
  \State $idx \gets y \cdot W + x$
  \If{$d > \texttt{dist}[idx]$}
    \State \textbf{continue}
  \EndIf
  \If{$(x,y) = t$}
    \State \textbf{break}
  \EndIf
  \ForAll{$(dx,dy) \in \Delta$}
    \State $nx \gets x + dx$, $ny \gets y + dy$
    \If{$0 \le nx < W$ and $0 \le ny < H$}
      \State $nidx \gets ny \cdot W + nx$
      \State $w \gets c(nx,ny)$
      \If{\texttt{dist[$idx$]} + $w < \texttt{dist[$nidx$]}$}
        \State \texttt{dist[$nidx$]} $\gets \texttt{dist[$idx$]} + w$
        \State \texttt{prev[$nidx$]} $\gets idx$
        \State Insert $(\texttt{dist[$nidx$]}, nx, ny)$ into $PQ$
      \EndIf
    \EndIf
  \EndFor
\EndWhile

\State Reconstruct path from $t$ using \texttt{prev} as before.
\end{algorithmic}
\end{algorithm}

\subsection{Enhancements and Advancement}

\begin{itemize}
  \item \textbf{Enhancement:} Introduces a cost function over pixels,
  allowing preference for thicker roads, brighter pixels, or specific
  classes of terrain.
  \item \textbf{Advancement:} Generalizes unweighted BFS to weighted
  shortest paths, enabling more realistic routing that accounts for
  road quality, width, or other image-derived attributes.
\end{itemize}
\clearpage
%%%%%%%%%%%%%%%%%%%%%%%%%%%%%%%%%%%%%%%%%%%%%%%%%%%%%%%%%%%%
\section{Left/Right/Up/Down Directional Solver}
\label{sec:lrud_solver}

\subsection{Concept}

The left/right/up/down directional solver is a constrained BFS that
only allows moves in the four cardinal directions, but can incorporate
directional penalties or preferences. For example, it can penalize
frequent turns or prefer forward motion, which is useful for generating
turn-minimizing routes.

\subsection{Pseudocode (BFS with Direction State)}

We augment the state with the current direction $dir \in \{L,R,U,D\}$.

\begin{algorithm}[h!]
\caption{Directional BFS with Turn Penalty}
\begin{algorithmic}[1]
\Require Start pixel $s$, target pixel $t$
\Ensure Path that balances distance and turn minimization

\State Directions:
\State \hspace{1em} $D = \{(1,0):R, (-1,0):L, (0,1):D, (0,-1):U\}$

\State For each pixel index $idx$ and direction $d$, maintain
\texttt{dist[idx][d]} (cost) and \texttt{prev[idx][d]}.

\State Initialize all \texttt{dist} to $+\infty$, \texttt{prev} to
\texttt{null}.

\State Use a priority queue $PQ$ over states $(cost,x,y,dir)$.

\State For all initial directions $d \in \{L,R,U,D\}$:
\State \hspace{1em} $idx_s \gets s_y \cdot W + s_x$
\State \hspace{1em} \texttt{dist[$idx_s$][d]} $\gets 0$
\State \hspace{1em} Insert $(0,s_x,s_y,d)$ into $PQ$

\State Turn penalty $\lambda > 0$.

\While{$PQ$ not empty}
  \State Extract $(c,x,y,dir)$ with minimal $c$
  \State $idx \gets y \cdot W + x$
  \If{$(x,y) = t$}
    \State \textbf{break}
  \EndIf
  \ForAll{$(dx,dy,dir2) \in D$}
    \State $nx \gets x + dx$, $ny \gets y + dy$
    \If{$0 \le nx < W$ and $0 \le ny < H$ and \texttt{isRoad(nx,ny)}}
      \State $nidx \gets ny \cdot W + nx$
      \State $turnCost \gets 0$ if $dir2 = dir$ else $\lambda$
      \State $newCost \gets c + 1 + turnCost$
      \If{$newCost < \texttt{dist[$nidx$][dir2]}$}
        \State \texttt{dist[$nidx$][dir2]} $\gets newCost$
        \State \texttt{prev[$nidx$][dir2]} $\gets (idx,dir)$
        \State Insert $(newCost,nx,ny,dir2)$ into $PQ$
      \EndIf
    \EndIf
  \EndFor
\EndWhile

\State Choose the best direction $d^\ast$ at $t$ with minimal
\texttt{dist[$idx_t$][d]} and reconstruct the path using \texttt{prev}.
\end{algorithmic}
\end{algorithm}

\subsection{Enhancements and Advancement}

\begin{itemize}
  \item \textbf{Enhancement:} Incorporates direction as part of the
  state and penalizes turns, allowing control over route smoothness and
  turn count.
  \item \textbf{Advancement:} Extends the basic 4-connected BFS by
  embedding directional preferences, bridging the gap between pure
  shortest-path routing and navigation-style routing that cares about
  turn minimization.
\end{itemize}
\clearpage
%%%%%%%%%%%%%%%%%%%%%%%%%%%%%%%%%%%%%%%%%%%%%%%%%%%%%%%%%%%%
\section{Overall Progression and Enhancements}

The algorithms form a progression:
\begin{enumerate}
  \item \textbf{Pixel-based maze solver:} Baseline unweighted BFS on a
  binary road image.
  \item \textbf{Longest-white-path solver:} Changes the objective from
  shortest to longest path, exploring coverage.
  \item \textbf{16-direction BFS solver:} Increases directional
  resolution, producing smoother paths.
  \item \textbf{Grid-based neighbor solver:} Introduces multi-resolution
  routing for scalability.
  \item \textbf{Image-based traversal solver:} Generalizes to weighted
  costs derived from image properties.
  \item \textbf{Left/right/up/down directional solver:} Adds direction
  to the state and penalizes turns, aligning with navigation semantics.
\end{enumerate}

Each step adds either:
\begin{itemize}
  \item A richer objective (longest path, turn minimization),
  \item A richer state (direction, cost),
  \item Or a richer representation (multi-resolution grid, weighted
  image).
\end{itemize}

\end{document}
